\section{Tổng kết dự án}

\subsection{Kết quả đạt được}
Sau thời gian tìm hiểu, nghiên cứu và triển khai, nhóm dự án đã hoàn thành các mục tiêu cơ bản được đề ra ban đầu. Cụ thể:
\begin{itemize}
    \item \textbf{Về mặt lý thuyết:} Đã tổng hợp và hệ thống hóa thành công các kiến thức liên quan đến IoT và Trí tuệ nhân tạo, tạo tiền đề vững chắc cho việc thiết kế và xây dựng hệ thống.
    \item \textbf{Về mặt thực tiễn:} Thiết kế và triển khai thành công Nhà thông minh áp dụng Trí tuệ nhân tạo.
    \item \textbf{Về mặt hiệu năng:} Hệ thống hoạt động tương đối ổn định trong môi trường thử nghiệm, giao diện trực quan, logic xử lý rõ ràng và thân thiện với người dùng.
\end{itemize}

\subsection{Hướng phát triển tương lai}
Nhằm hoàn thiện sản phẩm và nâng cao giá trị thực tiễn, định hướng phát triển của dự án trong thời gian tới bao gồm:
\begin{itemize}
    \item \textbf{Nâng cấp và mở rộng tính năng:} Nghiên cứu và tích hợp thêm Tiện ích nâng cao để hệ thống thông minh, tự động hóa và đáp ứng sát hơn nhu cầu của người dùng.
    \item \textbf{Tối ưu hóa hệ thống:} Áp dụng các kỹ thuật cấu trúc phần mềm hiện đại để cải thiện tốc độ phản hồi và khả năng mở rộng (scalability).
    \item \textbf{Thử nghiệm và triển khai thực tế:} Đưa sản phẩm vào thử nghiệm với tập người dùng đa dạng hơn để thu thập phản hồi (feedback), từ đó tinh chỉnh luồng nghiệp vụ và hoàn thiện các tiêu chuẩn bảo mật.
\end{itemize}